\section{Introdução}
\label{sec:introducao}

Dado o kernel do algoritmo sequencial para resolver o problema da Maior Subsequência Comum (\textit{Longest Common Subsequence} -- LCS) baseado em programação dinâmica clássica $\mathcal{O}(M \times N)$ (Código~\ref{lst:kernel_seq}), o objetivo deste trabalho é paralelizar esta computação utilizando os paradigmas de memória compartilhada (\textbf{OpenMP}), memória distribuída (\textbf{MPI}) e de forma \textbf{híbrida} (\textbf{MPI + OpenMP}), analisando o desempenho e os \textit{speedups} obtidos na arquitetura multicore da servidora \textit{duo} (AMD Ryzen 9 5950X, 16 núcleos físicos).

\begin{lstlisting}[caption={Kernel sequencial do algoritmo LCS em linguagem C.}, label={lst:kernel_seq}]
int LCS(mtype **scoreMatrix, int sizeA, int sizeB, char *seqA, char *seqB) {
    int i, j;
    for (i = 1; i < sizeB + 1; i++) {
        for (j = 1; j < sizeA + 1; j++) {
            if (seqA[j - 1] == seqB[i - 1]) {
                scoreMatrix[i][j] = scoreMatrix[i - 1][j - 1] + 1;
            } else {
                scoreMatrix[i][j] = max(scoreMatrix[i - 1][j], 
                                        scoreMatrix[i][j - 1]);
            }
        }
    }
    return scoreMatrix[sizeB][sizeA];
}
\end{lstlisting}
